\originalchapter{碰撞与散射}\label{ch:collision}
\originalsection{恢复系数与一维碰撞}
短时碰撞可首先利用动量守恒. 完全弹性碰撞还保持总动能; 完全非弹性碰撞后两物体粘合. 一般非弹性碰撞要给出额外模型, 不能仅凭动量唯一决定两个末速.

一维中两物体碰前速度 $u_1,u_2$, 碰后 $v_1,v_2$, 定义恢复系数 $e=-(v_2-v_1)/(u_2-u_1)$, 要求正在接近且碰后分离. 对通常耗散碰撞 $0\le e\le1$, 超弹性过程释放内部能可有 $e>1$. 合用总动量守恒,
\begin{align}
v_1&=\frac{(m_1-em_2)u_1+(1+e)m_2u_2}{m_1+m_2},\\
v_2&=\frac{(1+e)m_1u_1+(m_2-em_1)u_2}{m_1+m_2}.
\end{align}
将动能按式\eqref{eq:cmenergy}分解, 碰前后质心速度不变, 相对速度绝对值变为 $e$ 倍, 故动能损失为 $\frac12\mu(1-e^2)(u_1-u_2)^2$. $e=0$ 时损失最大; 并非全部初始动能都损失, 质心平动能仍保留.

\originalsection{质心系与二维弹性碰撞}
二维中动量守恒给两个分量方程, 再加弹性条件也通常不足以唯一指定两物体出射方向, 还需碰撞几何或相互作用模型. 在质心系中总动量为零, 两物体动量等大反向, 弹性碰撞只转动这对向量而保持大小. 这比在实验室系同时处理许多平方项更直观.

质量 $m_1$ 以速率 $u$ 打向静止的 $m_2$, 令质心散射角为 $\chi$. 质心速度沿初始方向, $V=m_1u/(m_1+m_2)$; 入射物在质心系的速率 $u'=m_2u/(m_1+m_2)$. 出射到实验室系后,
\begin{equation}
v_{1x}=V+u'\cos\chi,\quad v_{1y}=u'\sin\chi,\quad
\tan\theta_{\rm lab}=\frac{m_2\sin\chi}{m_1+m_2\cos\chi}.
\end{equation}
确定角象限需用两个分量, 不能只取反正切主值. 等质量时出射两速度垂直, 可从动量 $\vct u=\vct v_1+\vct v_2$ 与能量 $u^2=v_1^2+v_2^2$ 相减得到 $\vct v_1\cdot\vct v_2=0$. 单个速度为零的退化情形没有确定的出射方向.

\originalsection{弹道摆与多重碰撞}
子弹质量 $m$、速度 $u$, 射入静止摆块质量 $M$ 并粘住时, 碰撞阶段水平动量守恒给 $V=mu/(M+m)$. 随后摆升高 $h$ 时机械能守恒给 $V=\sqrt{2gh}$. 于是 $u=(M+m)\sqrt{2gh}/m$. 不能在子弹粘住的那一瞬间使用机械能守恒. 若摆块由刚性杆支持, 枢轴可能提供明显冲量, 则选枢轴角动量守恒更加稳妥.

多球碰撞中, 先找时间次序, 每次只对实际碰撞的物体使用守恒关系. 大球从地面反弹后碰小球, 在质量比趋于零的近似下大球像匀速移动的墙: 小球相对墙的速度反号. 在地面系有 $v'_{\rm small}=2v_{\rm wall}-v_{\rm small}$. 反复迭代可以放大上层球速度, 但假定无限弹性、很小质量比以及良好碰撞次序; 现实能量不能超过整体初始能量.
